\subsection{Feasibility of On-Device Training}

Recent advancements in optimization techniques have significantly reduced the memory and computational requirements for training machine learning models, making on-device training increasingly feasible, even on resource-constrained edge devices.
For instance, Lin et al. \cite{lin2022ondevice} demonstrated training neural networks on microcontroller units with only 256KB of RAM by employing an algorithm-system co-design framework. Their approach includes Quantization-Aware Scaling to stabilize 8-bit quantized training and Sparse Update to skip gradient computation of less important layers, making it particularly suitable for resource-constrained environments.
Similarly, Cai et al. \cite{cai2020tinytl} introduced a tiny transfer learning approach that freezes most parameters and only trains a small subset, allowing effective fine-tuning with minimal memory requirements. Qiu et al. \cite{qiu2022zerofl} proposed ZeroFL, a framework that relies on highly sparse operations to accelerate on-device training in federated learning settings, enabling efficient model training on edge devices with up to 95\% sparsity.
Recent work by Sugiura and Matsutani \cite{sugiura2025elasticzo} further advanced this field with ElasticZO, which combines zeroth-order and first-order optimization to achieve a better trade-off between accuracy and training cost. Their ElasticZO-INT8 variant achieves integer arithmetic-only training, further reducing memory usage and training time by approximately 1.5x without compromising accuracy.
% Zeroth-order methods have also been extended to federated learning for large language models as demonstrated by Ling et al. \cite{ling2024convergence}, who showed that memory-efficient zeroth-order optimization can reduce GPU memory usage during training to levels comparable to those during inference.
These advancements suggest that on-device training is no longer limited to high-end devices with abundant resources. Even small embedded systems with memory capacities measured in kilobytes rather than gigabytes can participate in model training.